\chapter{工程实现建议与部署方案}

\section{软件架构建议}
建议按“求解核心层、实验编排层、结果管理层”三层组织：
\begin{itemize}
  \item \textbf{核心层}：导数计算、积分器、优化器；
  \item \textbf{编排层}：参数扫描、批量试验、并行调度；
  \item \textbf{管理层}：日志、报表、异常追踪。
\end{itemize}

\section{在线计算流程}
在线场景可采用“离线标定 + 在线微调”：
\begin{enumerate}
  \item 离线建立先验参数与步长策略；
  \item 在线接收观测窗口，执行短窗迭代修正；
  \item 超阈值时触发重估与告警。
\end{enumerate}

\section{异常处理机制}
针对常见异常设置防护：
\begin{itemize}
  \item 雅可比病态：增大阻尼并切换线搜索；
  \item 误差突增：缩小步长并回滚上一步；
  \item 迭代超限：重置初值池并并行重启。
\end{itemize}

\section{部署性能建议}
在多核环境下将样本集分批并行，导数与残差评估可并行化。建议保存中间结果快照，便于故障恢复和结果复现。

\section{本章总结}
通过工程化封装，本文模型可从“单次实验代码”提升为“可持续运行组件”，满足批量试验、重复验证和在线校准需求。
